\chapter{预备知识}
\section{Plemelj 公式}
\begin{theorem}{Abel 定理} \label{ch1-thm-Abel}
    设 $ A \in \mathbb{C} $, 对矩阵微分方程 $ Y_{x} = AY $, 可得标量微分方程
    \begin{equation}
        (\det Y)_{x} = \tr(A) \det Y
    \end{equation}
    从而有
    \begin{equation}
        \det Y(x) = \det Y(x_{0}) \cdot e^{\int_{x_{0}}^{x} \rm{tr}(A(t)) \rmd t}
    \end{equation}
\end{theorem}
\begin{proof}
    由行列式的求导公式有,
    \begin{equation}
        (\det Y)_x = \tr\left(\adj(Y)Y_x\right),
    \end{equation}
    其中 $\adj(Y)$ 表示 $Y$ 的伴随矩阵. 代入 $Y_{x}=AY$, 并利用迹的循环不变性, 得
    \begin{equation}
        (\det Y)_{x} = \tr\left( \adj(Y) A Y \right) = \tr\left(Y\adj(Y)A\right).
    \end{equation}
    根据伴随矩阵恒等式 $Y \adj(Y)=(\det Y) I$ 可得
    \begin{equation}
        (\det Y)_{x} = \tr\left((\det Y) A\right) = \tr\bigl(A\bigr)\det Y
        \implies
        \det Y(x) = c_{0} e^{\int_{x_0}^{x}\tr\bigl(A(t)\bigr) \rmd t},
    \end{equation}
    再令 $x=x_{0}$, 有 $c_{0} = \det Y(x_0)$, 即证.
\end{proof}
\begin{theorem}{Morera 定理}\label{ch1-thm-Morera}
    如果 $ f(z) $ 在单连通区域 $ D $ 内连续, 且对于 $ D $ 内任意闭曲线 $ \Sigma $, 有 $ \int_{\Sigma} f(z) \rmd z = 0 $, 则有 $ f(z) $ 在 $ D $ 内解析
\end{theorem}
\begin{definition}{Schwartz 空间}
    欧几里得空间 $ \mathbb{R}^{n} $ 上的 Schwartz 空间 $ \mathcal{S}(\mathbb{R}^{n}) $ 定义为
    \begin{equation}
        \mathcal{S}(\mathbb{R}^{n}) = \left\{ f \in C^{\infty}(\mathbb{R^{n}}) : \| f \|_{\alpha, \beta} = \Sigma_{x \in \mathbb{R}^{n}} | x^{\alpha} \partial_{\beta}f(x) | < \infty, \alpha, \beta \in \mathbb{Z} \right\}
    \end{equation}
    其中 $ \alpha, \beta $ 为多重指标, $ \partial_{\beta} = \partial_{x_{1}}^{\beta_{1}} \cdots \partial_{x_{n}}^{\beta_{n}} $. $ f $ 称为速降函数或 Schwartz 函数.
\end{definition}
简单来说，速降函数是指当 $ x \to \pm \infty $ 时趋近于零的速度比所有的多项式的倒数都快，并且任意阶的导数都有这种性质的函数。
\begin{definition}[Bellmann-Gronwall 不等式]\label{ch1-def-Gronwall}
    设 $ u(t) $ 和 $ g(t) $ 是定义在区间 $ [t_{0}, T] $ 上的非负连续函数, 且满足
    \begin{equation}
        u(t) \leq A + \int_{t_{0}}^{t} g(s) u(s) \rmd s, \quad t \in [t_{0}, T]
    \end{equation}
    其中 $ A $ 为常数, 则
    \begin{equation}
        u(t) \leq A \exp{\int_{t_{0}}^{t} g(s) \rmd s}, \quad t \in [t_{0}, T]
    \end{equation}
\end{definition}
\begin{theorem}{Painleve 开拓定理}\label{ch1-thm-Painleve}
    设 $ D_{1}, D_{2} $ 为两个没有公共点的区域, 边界为 $ \Gamma $, 并设 $ f_{1}(z), f_{2}(z) $ 分别在 $ D_{1}, D_{2} $ 内解析, 在 $ D_{1} + \Gamma, D_{2} + \Gamma $ 上连续, 且 $ f_{1}(z) = f_{2}(z), \forall z \in \Gamma $, 则
    \begin{equation}
        f(z) = \begin{cases}
            f_{1}(z)            & z \in D_{1}  \\
            f_{1}(z) = f_{2}(z) & z \in \Gamma \\
            f_{2}(z)            & z \in D_{2}
        \end{cases}
    \end{equation}
    在 $ D_{1} + D_{2} + \Gamma $ 内解析.
\end{theorem}
\begin{proof}
    显然 $ f(Z) $ 在 $ D_{1} + D_{2} + \Gamma $ 上连续, 根据 Morera 定理, 只需证明 $ f(z) $ 沿任何闭曲线 $ \Sigma $ 的积分为 $ 0 $.

    如果 $ \Sigma \subset D_{1} $, 或 $ \Sigma \subset D_{2} $, 则由 $ f(z) $ 的解析性可得
    \begin{equation}
        \int_{\Sigma} f(z) \rmd z = 0
    \end{equation}
    如果 $ \Sigma $ 同时包含于 $ D_{1}, D_{2} $ 内, 把 $ \Gamma $ 在 $ \Sigma $ 内的曲线记为 $ C_{\gamma} $, 则
    \begin{equation}
        \int_{C_{1} + C_{\gamma}} f(z) \rmd z = 0, \quad \int_{C_{1} + C_{\gamma}}^{-} f(z) \rmd z = 0
    \end{equation}
    故
    \begin{equation}
        \int_{\Sigma} f(z) \rmd z =  \int_{C_{1} + C_{2}} f(z) \rmd z = \int_{C_{1} + C_{\gamma}} f(z) \rmd z = \int_{C_{1} + C_{\gamma}^{-}} f(z) \rmd z = 0.
    \end{equation}
\end{proof}
\begin{lemma}
    设  $ f(\xi) $ 在 $ z \in \Sigma$ 满足 $ \mu $ 次 H\"older 条件, 且 $ z' \to z$ 时, $ h/d $ 有界, 其中
    \begin{equation}
        h = | z' - z|, \quad d = \min_{\xi \in \Sigma} |\xi - z' | ,
    \end{equation} 则
    \begin{equation}
        \lim_{z' \to z} \frac{1}{2 \pi \rmi} \int_{\Sigma} \frac{f(\xi) - f(z)}{\xi - z'} \rmd \xi = \frac{1}{2 \pi \rmi} \int_{\Sigma} \frac{f(\xi) - f(z)}{\xi - z} \rmd \xi
    \end{equation}
\end{lemma}
\begin{proof}
    \begin{equation}
        \begin{aligned}
              & \frac{1}{2 \pi \rmi} \int_{\Sigma} \frac{f(\xi) - f(z)}{\xi - z'} \,\rmd \xi - \frac{1}{2 \pi \rmi} \int_{\Sigma} \frac{f(\xi) - f(z)}{\xi - z} \,\rmd \xi                                                                              \\
            = & \frac{1}{2\pi \rmi} \int_{\Sigma} \frac{(\xi - z)(f(\xi) - f(z)) -(\xi - z')(f(\xi) - f(z)) }{(\xi - z')(\xi - z)} \,\rmd \xi = \frac{1}{2\pi \rmi} \int_{\Sigma} \frac{z' - z}{\xi - z} \cdot \frac{f(\xi) - f(z)}{\xi - z} \,\rmd \xi \\
            = & \frac{1}{2\pi \rmi} \int_{\Sigma} \frac{h}{d} M \frac{(\xi - z)^{\mu}}{\xi - z} \, \rmd \xi = \Delta_{1} + \Delta_{2}
        \end{aligned}
    \end{equation}
    其中
    \begin{equation}
        \begin{aligned}
            | \Delta_{1} | & \leq \frac{hM}{2 \pi d} \int_{C_{\delta}} \frac{(\xi - z)^{\mu}}{\xi - z} \, \rmd | \xi | \leq \frac{hM}{2 \pi d} \int_{0}^{\delta} t^{\mu - 1} \, \rmd t = \frac{hM}{2 \pi d \delta} \delta^{\mu} (C_{\delta} = \{ | \xi - z | \leq \delta \} ) \\
            | \Delta_{2} | & =  \left| \frac{1}{2\pi} \int_{\Sigma \setminus  C_{\delta}} \frac{(f(\xi) - f(z))(z-z')}{\xi - z} \, \rmd \xi \right|
        \end{aligned}
    \end{equation}
    由于 $ \Sigma \setminus  C_{\delta} $ 不包含 $ z $, 故 $ \Delta_2 $ 为关于 $ z' $ 的连续函数, 则
    \begin{equation}
        | \Delta_{2} | =  \left| \frac{1}{2\pi \rmi} \int_{\Sigma \setminus  C_{\delta}} \frac{f(\xi) - f(z)}{\xi - z'} \, \rmd \xi - \frac{1}{2\pi} \int_{\Sigma \setminus  C_{\delta}} \frac{f(\xi) - f(z)}{\xi - z} \, \rmd \xi \right| < | \Delta_1 |
    \end{equation}
    取 $ \delta \to 0 $, 可令 $ | \Delta_1 | < \frac{\epsilon}{2} $, 故 $ | \Delta_1 | + | \Delta_2 | < \epsilon $
\end{proof}

\begin{theorem}{Plemelj 定理}
    设 $ z \in \Sigma $ 为正则点, 且不为边界点, $ f(\xi) $ 在 $ z $ 点满足 $ \mu $ 次 H\"older 条件, 且 $ z' \to z $ 时, $ h/d $ 有界, 其中 $ h = | z' - z|, \quad d = \min_{\xi \in \Sigma} |\xi - z' | $, 则
    \begin{align}
        F_{+} & = \lim_{z' \to z} F(z') = F(z) + \frac{1}{2}f(z), \quad z' \in \Sigma_{+}, \label{ch1-eq-Plemelj-1} \\
        F_{-} & = \lim_{z' \to z} F(z') = F(z) - \frac{1}{2}f(z), \quad z' \in \Sigma_{-}, \label{ch1-eq-Plemelj-2}
    \end{align}
\end{theorem}
\begin{proof}
    首先证明闭区线情形: $\forall z' \in \Sigma_{+}$
    \begin{equation}
        \begin{aligned}
            F_{+} & = \lim_{z' \to z} F(z') = \lim_{z' - z} \int_{\Sigma} \frac{1}{2 \pi \rmi} \frac{f(\xi)}{\xi - z'} \, \rmd \xi                                                            \\
                  & = \lim_{z' - z} \int_{\Sigma} \frac{1}{2 \pi \rmi} \frac{f(\xi) - f(z)}{\xi - z'} \, \rmd \xi + \frac{f(\xi)}{2 \pi \rmi} \int_{\Sigma} \frac{1}{\xi - z'} \, \rmd \xi    \\
                  & = \lim_{z' - z} \frac{1}{2 \pi \rmi} \int_{\Sigma_{+}}  \frac{f(\xi) - f(z)}{\xi - z'}  \rmd \xi + f(z)                                                                   \\
                  & = \lim_{z' - z} \frac{1}{2 \pi \rmi} \int_{\Sigma_{+}}  \frac{f(\xi)}{\xi - z'}  \rmd \xi - \frac{f(z)}{2 \pi \rmi} \int_{\Sigma_{+}}  \frac{1}{\xi - z'}  \rmd \xi+ f(z) \\
                  & = F(z) + \frac{1}{2}f(z)
        \end{aligned}
    \end{equation}
    $ F_{-}(z') $ 同理, 下证 $ \Sigma $ 为开曲线情形, 则可补充 $ \Sigma' $, s.t. $ \Sigma \cup \Sigma' $ 为闭曲线, 且定义 $ \forall \xi \in \Sigma', f(\xi) = 0 $, 则
    \begin{equation}
        F(z) = \frac{1}{2 \pi \rmi} \int_{\Sigma}  \frac{f(z)}{\xi - z}  \rmd \xi = \frac{1}{2 \pi \rmi} \int_{\Sigma \cup \Sigma'}  \frac{f(z)}{\xi - z}  \rmd \xi.
    \end{equation}
    则由开曲线情形可得
    \begin{equation}
        \begin{aligned}
            F_{+}(z) & = \lim_{\substack{z' \to z                                                                                                                                                      \\ z' \in +\Sigma}} \int_{\Sigma \cup \Sigma'}  \frac{f(\xi)}{\xi - z}  \rmd \xi = \int_{\Sigma \cup \Sigma'}  \frac{f(\xi) - f(z)}{\xi - z}  \rmd \xi + \frac{f(z)}{2 \pi \rmi} \int_{\Sigma \cup \Sigma'} \frac{1}{\xi - z'} \rmd \xi \\
                     & = \frac{1}{2 \pi \rmi} \int_{\Sigma} \frac{f(\xi) - f(z)}{\xi - z} \rmd \xi - \frac{f(z)}{2 \pi \rmi} \int_{\Sigma'} \frac{1}{\xi - z} \rmd \xi + f(z) = F(z) + \frac{1}{2}f(z)
        \end{aligned}
    \end{equation}
    $ F_{-}(z') $ 同理
\end{proof}

\begin{remark}
    如果 $ z \in \Sigma $ 为一个角点, 在其两切线的夹角为 $ \alpha $, 则可
    \begin{equation}\begin{aligned}
            F_{+}(z) & = \lim_{\substack{z' \to z \newline z' \in D}} \frac{1}{2 \pi \rmi} \int_{\Sigma}  \frac{f(\xi)}{\xi - z}  \rmd \xi = \frac{1}{2 \pi \rmi} \int_{\Sigma}  \frac{f(\xi) - f(z)}{\xi - z}  \rmd \xi + \frac{\alpha}{2 \pi} f(z) + (1 - \frac{\alpha}{2 \pi}) f(z) \\
                     & = F(z) - \frac{\alpha}{2 \pi} f(z) \label{ch1-eq-arc-Plemelj-1}
        \end{aligned}\end{equation}
    \begin{equation}\begin{aligned}
            F_{-}(z) & = \frac{1}{2 \pi \rmi} \int_{\Sigma}  \frac{f(\xi) - f(z)}{\xi - z}  \rmd \xi = \frac{1}{2 \pi \rmi} \int_{\Sigma}  \frac{f(\xi) - f(z)}{\xi - z}  \rmd \xi + \frac{\alpha}{2 \pi} f(z) - \frac{\alpha}{2 \pi} f(z) \\
                     & = F(z) - \frac{\alpha}{2 \pi} f(z) \label{ch1-eq-arc-Plemelj-2}
        \end{aligned}\end{equation}
\end{remark}
\begin{definition}{Plemelj 公式}
    由 \eqref{ch1-eq-Plemelj-1} - \eqref{ch1-eq-Plemelj-2}, \eqref{ch1-eq-arc-Plemelj-1} - \eqref{ch1-eq-arc-Plemelj-2} 可以看出, 无论 $ z $ 是正则点还是角点, 都有
    \begin{align}
        F_{+}(z) - F_{-}(z) & = f(z), \quad z \in \Sigma                                                    \\
        F_{+}(z) - F_{-}(z) & = \frac{1}{\pi \rmi} \int_{\Sigma} \frac{f(\xi)}{\xi - z} \rmd \xi =: H(f)(z)
    \end{align}
    称为标量 RH 问题, 其解可用 Cauchy 积分给出
    \begin{equation}
        F(z) = A + \frac{1}{2 \pi \rmi} \int_{\Sigma} \frac{f(\xi)}{\xi - z} \rmd \xi, \quad z \in \mathbb{C}
    \end{equation}
    其中 $ A $ 为任意常数, 一般被边值或渐进条件决定, 这一公式被称为 Plemelj 公式.
\end{definition}


对于如下 RH 问题
\begin{equation}
    G_{+}(z) = G_{-}(z)v(z), \quad z\in \Sigma
\end{equation}
只需两边取对数变换则有
\begin{equation}
    \ln{G_{+}(z)}- \ln{G_{-}(z)} =  \ln{v(z)}, \quad z\in \Sigma
\end{equation}
由 Plemelj 公式, 则有
\begin{equation}
    \ln{G(z)} = A + \frac{1}{2 \pi \rmi} \int_{\Sigma} \frac{v(\xi)}{\xi - z} \rmd \xi \implies G(z) = B \exp{\left( \frac{1}{2 \pi \rmi} \int_{\Sigma} \frac{f(\xi)}{\xi - z} \rmd \xi \right)}, \quad z\in \Sigma
\end{equation}
其中 $ B $ 为任意常数.

\section{矩阵 RH 问题}
\begin{definition}{RH 问题}
    设 $ \Sigma $ 为复平面 $ \mathbb{C} $ 内的有向路径, 假设存在一个 $ {\Sigma} ^{0} $ 上的光滑映射 $ v(z) : \Sigma \to GL(n,\mathbb{C}) $, 则 $ (\Sigma, v) $ 决定了一个 RH 问题, 寻找一个 $ n $ 阶矩阵 $ M(z) $ 满足
    \begin{align}
        M(z) \in C, \quad ( \mathbb{C}\setminus  \Sigma ) \label{ch1-eq-RH-analytical-condition} \\
        M_{+}(z) = M_{-}(z)v(z), \quad z\in \Sigma \label{ch1-eq-RH-Matrix-jump-condition}       \\
        M(z) \to I ,\quad z \to \infty \label{ch1-eq-RH-Matrix-asymptotic-condition}
    \end{align}
    其中 $ M_{\pm} $ 表示在正负区域内 $ z' \to z $ 时的极限, $ \Sigma $ 称为跳跃曲线, $ v(z) $ 称为跳跃矩阵.
\end{definition}

根据 Beadls-Coifman 定理, 如上 RH 问题的解可通过如下方式构造: 不妨设跳跃矩阵 $ v(z) $ 具有如下分解
\begin{equation}
    v = (b_{-})^{-1}b_{+}
\end{equation}
由此, 可构造
\begin{equation}
    w_{+} = b_{+} - I, \quad w_{-} = I - b_{-}
\end{equation}
进一步可定义 Cauchy 投影算子
\begin{equation}
    (C_{\pm} f)(z) = \lim_{\substack{z' \to z \in \Sigma \\ z' \in \pm \Sigma}} \frac{1}{2 \pi \rmi} \int_{\Sigma} \frac{f(\xi)}{\xi - z'} \rmd \xi \label{ch1-eq-Cauchy-projection}
\end{equation}
则可以证明如果 $ f(z) \in L^{2}(\Sigma) $, 则 $ C_{\pm} : L^{2} \to L^{2} $ 的有界算子, 且 $ C_{+} - C_{-} = 1 $, 再定义算子
\begin{equation}
    C_{w}f = C_{+}(fw_{-}) - C_{-}(fw_{+}) \label{ch1-eq-Cauchy-projection-operator}
\end{equation}
则 $ C_{w}: L^{2} \cap L^{\infty} \to L^2 $ 的有界算子.

\begin{theorem}{Beals--Coifman 定理}
    设 $ \det v = 1 $, 算子 $ I - C_{w} $ 在 $ L^{2}(\Sigma) $ 上可逆, $ \mu \in I +L^{2}(\Sigma) $ 为下列方程
    \begin{equation}
        (I - C_{w}) \mu = I \label{ch1-eq-RH-singular-integral}
    \end{equation}
    的解, 且
    \begin{equation}
        (I -C_{w})(\mu - I) = C_{w}I + C_{+}w_{-} + C_{-}w_{+} \in L^{2}(\Sigma)
    \end{equation}
    则
    \begin{equation}
        M(z) = I + \frac{1}{2 \pi \rmi} \int_{\Sigma} \frac{\mu(\xi) w(\xi)}{\xi - z} \rmd \xi \label{ch1-eq-RH-singular-integral-matrix}
    \end{equation}
    为上述 RH 问题的唯一解, RH 问题 $ M(z) $ 的可解等价于奇异积分方程 \eqref{ch1-eq-RH-singular-integral}.
\end{theorem}
\begin{proof}
    只需证明 \eqref{ch1-eq-RH-singular-integral-matrix} 满足 \eqref{ch1-eq-RH-singular-integral}
    \begin{equation}
        \begin{aligned}
            M_{+} & = I + \lim_{z' \to z \in \Sigma} \frac{1}{2 \pi \rmi }\int_{\Sigma} \frac{\mu(\xi) w(\xi)}{\xi - z} \rmd \xi \overset{\eqref{ch1-eq-Cauchy-projection}}{=} I + C_{+}(\mu w)      \\
                  & = I + C_{+}(\mu w_{-}) + C_{+}(\mu w_{+}) = I + C_{+}(\mu w_{-}) + C_{-}(\mu w_{+}) +   C_{+}(\mu w_{+}) + C_{-}(\mu w_{+})                                                      \\
                  & \overset{\eqref{ch1-eq-Cauchy-projection-operator}}{=} I + C_{w}(\mu) + \mu w_{+} \overset{\eqref{ch1-eq-RH-singular-integral}}{=} \mu + \mu w_{+} = \mu (I + w_{+}) = \mu b_{+}
        \end{aligned}
    \end{equation}
    同理有
    \begin{equation}
        M_{-}(z) = \mu b_{-}
    \end{equation}
    故有
    \begin{equation}
        M_{+}(z) = \mu b_{+} = \mu b_{-}(b_{-})^{-1} b_{+} = M_{-}(z)v(z)
    \end{equation}
    下证唯一性, 先证 $ M $ 可逆, 对 \eqref{ch1-eq-RH-Matrix-jump-condition} 取行列式, 并注意到 $ \det v(z) = 1 $, 则 $ \det M_{+}(z) = \det M_{-}(z) $. 故由 Painleve 开拓定理得 $ \det M(z) $ 在 $ \mathbb{C} $ 上解析. 由 \eqref{ch1-eq-RH-Matrix-asymptotic-condition} 得
    \begin{equation}
        \det M(z) \to 1, z \to \infty
    \end{equation}
    故由 Liouville 定理可知 $ \det M(z) = c $, 再由渐进条件得 $ c = 1 $, 故 $ M(z) $ 可逆.

    设 $ \widetilde{M} $ 为上述 RH 问题的另一个解, 则 $ \widetilde{M} $ 可逆, 且在 $ \mathbb{C} \setminus  \Sigma $ 上解析, 且在 $ \Sigma $ 上满足
    \begin{equation}
        (M \widetilde{M}^{-1}) = M_{+} \widetilde{M}_{+}^{-1} = M_{-} v (\widetilde{M}_{-} v)^{-1} = M_{-}(\widetilde{M}_{-})^{-1} = (M \widetilde{M}^{-1})_{-}
    \end{equation}
    由 Painleve 开拓定理, $ M \widetilde{M}^{-1} $ 在 $ \mathbb{C} $ 上解析. 另外由 $ M, \widetilde{M} \to I $, 故 $ M \widetilde{M}^{-1} $ 有界, 故为常矩阵. 由渐进性得
    \begin{equation}
        M \widetilde{M}^{-1} = I \implies M = \widetilde{M}
    \end{equation}
\end{proof}

\begin{remark}
    由 RH 问题解的唯一性, 解与跳跃矩阵 $ v(z) $ 的分解无关, 因此可以考虑 $ v(z) $ 的平凡解
    \begin{equation}
        b_{-} = I, \quad b_{+} = v,
    \end{equation}
    故可构造 $w_{-} = 0, \quad w_{+} = v - I, \quad w = v - I$, 则
    \begin{equation}
        C_{w}f =  C_{+}(fw_{-}) + C_{-}(fw_{+}) = C_{+}(f(v - I)), \quad \mu = (I - C_{w})^{-1}I
    \end{equation}
    则 RH 问题的解可表示为
    \begin{equation}
        M(z) = I + \frac{1}{2 \pi \rmi} \int_{+\Sigma} \frac{\mu(\xi) (v(\xi) - I)}{\xi - z}  \rmd \xi \label{ch1-eq-RH-solution}
    \end{equation}
\end{remark}

RH 方法的关键思想就是改变积分路径, 通过跳跃矩阵的分解情况, 确定对积分路径进行一系列形变, 再取极限去除跳跃矩阵为单位阵的情形, 将其化解为可解的 RH 问题, 所以一个自然的问题就是:「为什么可以扔掉跳跃矩阵为单位阵的路径? 或者说为什么跳跃矩阵对对 RH 问题的解不产生贡献.」

这很容易通过表达式 \eqref{ch1-eq-RH-solution} 看出, 我们将积分路径分解为 $ \Sigma = \Sigma_{1} + \Sigma_{2} $ 且
\begin{equation}
    v = \begin{cases}
        v_1 \neq I \quad \Sigma_{1}, \\
        v_1 = I \quad \Sigma_{2}
    \end{cases}
\end{equation}
则在 $ \Sigma_{1} $ 上, $ v- I = v_{1} - I \neq 0 $, 在 $ \Sigma_{2} $ 上, $ v- I = v_{1} - I = 0 $, 故
\begin{equation}
    \begin{aligned}
        M(z) & = I + \frac{1}{2 \pi \rmi} \int_{\Sigma} \frac{\mu(\xi) (v(\xi) - I)}{\xi - z}  \rmd \xi                                                                                                      \\
             & = I + \frac{1}{2 \pi \rmi} \int_{\Sigma_{1}} \frac{\mu(\xi) (v_{1}(\xi) - I)}{\xi - z}  \rmd \xi + \frac{1}{2 \pi \rmi} \int_{\Sigma_{2}} \frac{\mu(\xi) (v_{2}(\xi) - I)}{\xi - z}  \rmd \xi \\
             & = I + \frac{1}{2 \pi \rmi} \int_{\Sigma_1} \frac{\mu(\xi) (v_{1}(\xi) - I)}{\xi - z}  \rmd \xi
    \end{aligned}
\end{equation}
易得 RH 问题 $ (\Sigma, v(z)) $ 的解与 $ (\Sigma_{1}, v_{1}(z)) $ 的解相同, 即跳跃矩阵为单位阵的路径 $ \Sigma_{2} $ 对 RH 问题的解无贡献, 可以舍去. 只求 $ \Sigma_{1} $ 上的 RH 问题

\begin{lemma}\label{ch1-lem-commuting-matrices}
    设 $A, B$ 为 $N$ 阶复矩阵, 且 $AB=BA$, 对任意特征值 $\lambda$, 矩阵 $A$ 的特征子空间 $E_{\lambda}(A) := \ker(A-\lambda I)$ 在矩阵 $B$ 的作用下保持不变, 即
    \begin{equation}
        B E_{\lambda}(A) \subseteq E_{\lambda}(A).
    \end{equation}
    特别地，若 $A$ 有 $n$ 个互不相同的特征值
    $\lambda_{1},\ldots,\lambda_{n}$，且
    \begin{equation}
        A r_{i} =\lambda_{i} r_{i}, \quad i = 1, \ldots, n,
    \end{equation}
    则存在 $\mu_i\in\mathbb{C}$ 使得
    \begin{equation}
        B r_{i}=\mu_{i} r_{i}, \quad i = 1, \ldots, n.
    \end{equation}
    因而 $r_1,\ldots,r_n$ 构成 $A,B$ 的一组公共特征向量基, 且 $A,B$ 可同时对角化。
\end{lemma}

\begin{proof}
    任取 $v\in E_{\lambda}(A)$, 则
    \begin{equation}
        (A-\lambda I)Bv = B(A-\lambda I)v = 0.
    \end{equation}
    因而 $Bv\in E_{\lambda}(A)$, 即
    $B E_{\lambda}(A)\subseteq E_{\lambda}(A)$。

    设 $A$ 有 $n$ 个互异特征值, $E_{\lambda_{i}}(A) = \operatorname{span} \{r_i\}$. 由于 $B r_{i} \in E_{\lambda_{i}}(A)$, $ \exists \mu_{i} \in \mathbb{C}$ 使得 $B r_{i}=\mu_{i} r_{i}$, 令 $R=(r_{1}, \ldots, r_{n})$, 则 $R$ 可逆, 且
    \begin{equation}
        R^{-1}AR = \diag(\lambda_1,\ldots,\lambda_n), \quad
        R^{-1}BR = \diag(\mu_1,\ldots,\mu_n).
    \end{equation}
\end{proof}
\begin{remark}
    由 $A B = B A$ 也易见 $B \rme^{A x}=\rme^{A x} B$.
\end{remark}
考虑常系数齐次线性方程组
\begin{equation}
    \frac{\rmd y}{\rmd x}
    =
    Ay.
    \label{ch1-ch1-eq-A-linear-system}
\end{equation}

\begin{corollary}
    设 $A\in\mathbb{C}^{n\times n}$ 有 $n$ 个互不相同的特征值 $\lambda_1,\ldots,\lambda_n$, 相应的特征向量为 $r_1, \ldots, r_n$. 则
    \begin{equation}
        \Phi(x) = \left( \rme^{\lambda_1x}r_1, \rme^{\lambda_2x}r_2, \ldots, \rme^{\lambda_nx}r_n \right) = R \diag\left( \rme^{\lambda_1x}, \ldots, \rme^{\lambda_nx} \right)
    \end{equation}
    为方程 \eqref{ch1-ch1-eq-A-linear-system} 的基解矩阵, 其中
    $R=(r_{1},\ldots,r_{n})$.
\end{corollary}
现考虑常系数偏微分方程组
\begin{equation}
    \frac{\partial y}{\partial x} = A y, \quad
    \frac{\partial y}{\partial t} = B y.
    \label{ch1-ch1-eq-AB-linear-system}
\end{equation}

\begin{corollary}
    设 $A, B$ 为 $n$ 阶矩阵, 满足 $AB=BA$, 且 $A$ 有 $n$ 个互不相同的特征值。记
    \begin{equation}
        A r_{i}=\lambda_{i} r_{i}, \quad B r_{i}=\mu_{i} r_{i}.
    \end{equation}
    则上述方程组的规范化基解矩阵为
    \begin{equation}
        F(x,t) = \rme^{Ax}\rme^{Bt} = \rme^{Ax+Bt} = R\diag\left( \rme^{\lambda_{1} x + \mu_{1} t}, \ldots, \rme^{\lambda_{n} x + \mu_{n} t} \right) R^{-1}.
    \end{equation}
    因而方程组的一般解为 $y(x,t) = \rme^{Ax+Bt}c_0$, 其中 $c_0\in\mathbb{C}^n$ 为常向量. 同理有
    \begin{equation}
        \Phi(x,t) = F(x,t) R = \left( e^{\lambda_{1} x r_{1}}, \ldots, e^{\lambda_{n} x r_{n}} \right)
    \end{equation}
    也是该方程组的基解矩阵。
\end{corollary}
\begin{proof}
    由 \eqref{ch1-ch1-eq-AB-linear-system}, 不妨设 $y(x,t)=\rme^{A x} c(t)$, 将其代入 \eqref{ch1-ch1-eq-AB-linear-system} 右式可得
    \begin{equation}
        \rme^{A x} c'(t) = B \rme^{A x}c(t) = \rme^{A x} B c(t) \implies c'(t)=B c(t) \implies c(t)=\rme^{Bt}c_{0},
    \end{equation}
    则 $y = \rme^{A x} \rme^{B t} c_{0} = \rme^{A x + B t} c_{0}$.
\end{proof}
